\documentclass[10pt,a4paper]{amsart}
\usepackage[margin=19mm]{geometry}
\usepackage{amsmath,amssymb,mathtools}
\usepackage[scr=rsfs,cal=euler]{mathalfa}
\usepackage{xfrac}
\usepackage[unicode,hidelinks,bookmarks=false]{hyperref}
\newcommand{\ie}{i.e.\ }
\newcommand{\cf}{cf.\ }
\newcommand{\resp}{respectively\ }
\newcommand{\Subsection}{\subsection}
\providecommand{\nobreakdash}{\nobreak\hbox{-}}
\setlength{\emergencystretch}{2em}
\multlinegap0pt
\usepackage{bm}
\usepackage{bbm}
\usepackage{dsfont}
\usepackage{xspace}
\usepackage{cancel}
\usepackage{mathtools}
\usepackage[shortlabels]{enumitem}
\usepackage{amsthm}
\makeatletter
\@ifundefined{theorem}{\newtheorem{theorem}{Theorem}}{}
\@ifundefined{lemma}{\newtheorem{lemma}[theorem]{Lemma}}{}
\@ifundefined{proposition}{\newtheorem{proposition}[theorem]{Proposition}}{}
\makeatother
\title[Riesz transform and its related inequalities for degenerate ]{Riesz transform and its related inequalities for degenerate elliptic operators of Grushin type}
\author{Dangyang He}
\date{Source: arXiv:2607.15599v1 (CC BY 4.0)}
\begin{document}
\maketitle

\noindent\emph{Source and licence.} Riesz transform and its related inequalities for degenerate elliptic operators of Grushin type. Version arXiv:2607.15599v1 (CC BY 4.0), distributed under the Creative Commons Attribution 4.0 International licence, as recorded in the arXiv licence metadata for this exact version and shown on the abstract page https://arxiv.org/abs/2607.15599v1. Full rights notice: licenses/arxiv-2607.15599-CC-BY-4.0.txt.\par\smallskip

\noindent\textit{Editorial scope.} Counted unit: the Proposition whose source label is \texttt{prop:Green:Shur} in the article (source0.tex lines 2318--2340, source label \texttt{prop:Green:Shur}), delivered together with the article's own complete original proof of it (source0.tex lines 2342--2797, one \texttt{proof} environment of 456 lines). No mathematics is written, repaired or re-derived: statement and proof are the article's own text line for line. Required context retained uncounted: lem\_bessel\_inputs\_high\_frequency (lines 1267--1289); eq\_S\_lower\_large (lines 1269--1274); eq\_K\_normalized\_exp\_decay (lines 1276--1282); eq\_kernel\_Green (lines 1706--1712); lem:Bessel:integral (lines 2268--2280). Every definition, display and label of those lines is retained unchanged. Omitted from this package and not redistributed: 1--1266, 1275--1275, 1283--1705, 1713--2267, 2281--2317, 2341--2341, 2798--6554. Nothing inside a retained excerpt is omitted or altered: every retained body is delivered exactly as the article's lines are written, so no omission receipt of any kind accompanies this package. Citations: the retained excerpts carry no citation command at all, so no bibliography entry of the article is transcribed and no reference is redistributed. Reference adaptation: none was needed and none was made; every reference inside the delivered unit resolves inside it, so no unresolved reference or citation is produced. Printed slips kept literal: no printed word, symbol, hypothesis or number of the counted statement or of its proof was altered. Layout and class: the article's own class is \texttt{amsart} with packages including amsmath, amsthm, amsfonts, amssymb, color, mathrsfs, enumitem, amstext, amsxtra, txfonts, hyperref, pgf, tikz, caption; the wrapper is the lane's pinned \texttt{amsart} editorial wrapper at 10pt on A4, which restates the theorem-like environments the retained text needs and adds the article's own macro definitions that the retained text needs (none), with extra wrapper packages (bm, bbm, dsfont, xspace, cancel, mathtools, enumitem). The delivered numbering is the unit's own. Assets: no retained span contains a figure environment, a graphics-inclusion command, a PDF-page inclusion, a table or a TikZ picture; no image file is copied into this package, the package declares no asset dependency and no image file is redistributed with it. Licence: version arXiv:2607.15599v1 of this article is distributed under the Creative Commons Attribution 4.0 International licence, as recorded in the arXiv licence metadata for this exact version and shown on the abstract page https://arxiv.org/abs/2607.15599v1. A full rights notice is delivered at licenses/arxiv-2607.15599-CC-BY-4.0.txt. Characters: every retained excerpt is pure ASCII, so the delivered engine, XeLaTeX with its default Latin Modern text font, reports no missing character.
\begin{lemma}\label{lem_bessel_inputs_high_frequency}
Let $\nu_*>-1$. There exist constants $c,C>0$, depending only on $\nu_*$, such that, for every $\nu\geq\nu_*$,
\begin{equation}\label{eq_S_lower_large}
    \frac{I_{\nu+1}(t)}{I_\nu(t)}
    \geq c,
    \qquad
    t\geq C(1+\nu).
\end{equation}
Moreover, the function $\mathcal K_\nu(t)=t^\nu K_\nu(t)$ is decreasing on $(0,\infty)$, and
\begin{equation}\label{eq_K_normalized_exp_decay}
    \frac{\mathcal K_\nu(t_2)}{\mathcal K_\nu(t_1)}
    \leq
    C\exp\bigl\{-c(t_2-t_1)\bigr\},
    \qquad
    t_2\geq t_1\geq C(1+\nu).
\end{equation}
In addition, for every $\nu\geq 0$ and $t>0$,
\begin{equation}\label{eq:product:Bessel}
    I_\nu(t)K_\nu(t)
    \leq
    \frac{C}{t+\nu}.
\end{equation}
\end{lemma}

\begin{align}\label{eq_kernel_Green}
    G_{\ell,Y}(\rho,\sigma) = \mathfrak a^{-1}
    \begin{cases}
        \phi_{\ell}(\rho,Y) \, \psi_{\ell}(\sigma,Y), & 0 < \rho < \sigma,\\[4pt]
        \phi_{\ell}(\sigma,Y) \, \psi_{\ell}(\rho,Y), & \sigma < \rho < 2.
    \end{cases}
\end{align}

\begin{lemma}\label{lem:Bessel:integral}
Let $\nu>-1$ and $s>|\nu|$. Let $\varepsilon \in \{0,1\}$. The following two estimates hold:
\begin{enumerate}[label=(\roman*)]
    \item if $q\ge -\nu - \varepsilon$ and $q+s\le 2$, then
    \begin{equation*}
        \sup_{t>0}I(q,s,\nu,\varepsilon,t):=\sup_{t>0} t^q I_{\nu+\varepsilon}(t) \int_t^\infty u^{s-1} K_{|\nu|}(u) du \le C,
    \end{equation*}
    \item if $q+s+\nu-|\nu|\ge 0$ and $q+s\le 2$, then
    \begin{equation*}
        \sup_{t>0}II(q,s,\nu,t):= \sup_{t>0} t^q K_{|\nu|}(t) \int_0^t u^{s-1} I_{\nu}(u) du \le C.
    \end{equation*}
\end{enumerate}
\end{lemma}

\begin{proposition}\label{prop:Green:Shur}
For $\ell\ge 1$ and $0<\rho\le 2$, we have
\begin{equation}\label{eq_basic_R_schur_improved}
    \int_0^2 G_{\ell,Y}(\rho,\sigma) \, \sigma^{n-1} \, d\sigma \le \begin{cases}
        C \rho^{1-\alpha} (1+\ell)^{-2}, & |Y|\le C^*(1+\ell),\\[4pt]
        C \rho^{1-\alpha} (1+\ell+\tau_\rho)^{-2}, & |Y|> C^*(1+\ell).
    \end{cases}
\end{equation}
In particular,
\begin{equation}\label{eq_basic_R_schur_ell=0}
    \int_0^2 G_{0,Y}(\rho,\sigma) \, \sigma^{n-1} \, d\sigma \le \begin{cases}
        C, & |Y|\le C^*,\\[4pt]
        C (1+\tau_\rho)^{-2}, & |Y|> C^*.
    \end{cases}
\end{equation}
Moreover, for $\ell \ge 0$,
\begin{align}\label{eq_basic_R_schur}
\sup_{0<\rho<2} \rho^{2\beta} \int_0^2 G_{\ell,Y}(\rho,\sigma) \, \sigma^{n-1} \, d\sigma \le \begin{cases}
        C (1+\ell)^{-2}, & |Y|\le C^*(1+\ell),\\[4pt]
        C |Y|^{-2}, & |Y|> C^*(1+\ell).
    \end{cases}
\end{align}
\end{proposition}

\begin{proof}
\textit{Case 1: $|Y|\le C^*(1+\ell)$.} 
%The finite-energy solution at $\rho=0$ is
% \begin{equation*}
%     \phi_{\ell,Y}(\rho) = \rho^{-\mathfrak b}I_{\nu_\ell}(\tau_\rho).
% \end{equation*}
% The solution satisfying the Dirichlet condition at $\rho=2$ is
% \begin{equation*}
%     \psi_{\ell,Y}(\rho) = \rho^{-\mathfrak b} \bigl[ K_{\nu_\ell}(\tau_\rho) - \eta_{\ell,Y} I_{\nu_\ell}(\tau_\rho)\bigr],\qquad \eta_{\ell,Y}
%     :=
%     \frac{K_{\nu_\ell}(\tau_2)}{I_{\nu_\ell}(\tau_2)}.
% \end{equation*}
The one-dimensional Green kernel takes the form (cf.~\eqref{eq_kernel_Green}):
\begin{align}\label{eq_kernel_Green_copy}
G_{\ell,Y}(\rho,\sigma) = \mathfrak a^{-1}
\begin{cases}
\phi_{\ell}(\rho,Y)\,\psi_{\ell}(\sigma,Y), &0<\rho<\sigma,\\[4pt]
\phi_{\ell}(\sigma,Y)\,\psi_{\ell}(\rho,Y), &\sigma < \rho<2,
\end{cases}
\end{align}
where $\phi_{\ell,Y}(\rho) = \rho^{-\mathfrak b}I_{\nu_\ell}(\tau_\rho)$ and
\begin{equation*}
\psi_{\ell,Y}(\rho) = \rho^{-\mathfrak b} \bigl[ K_{\nu_\ell}(\tau_\rho) - \eta_{\ell,Y} I_{\nu_\ell}(\tau_\rho)\bigr],\qquad \eta_{\ell,Y} := \frac{K_{\nu_\ell}(\tau_2)}{I_{\nu_\ell}(\tau_2)}.
\end{equation*}


Suppose first $\ell\ge 1$. With the notation
\begin{equation*}
\gamma_\ell:= -\mathfrak b+\mathfrak a\nu_\ell = -\mathfrak b+\sqrt{\mathfrak b^2+\lambda_\ell},\quad \delta_\ell := \mathfrak b+\mathfrak a\nu_\ell = \mathfrak b+\sqrt{\mathfrak b^2+\lambda_\ell},
\end{equation*}
both $\gamma_\ell$ and $\delta_\ell$ are of size $1+\ell$. Since $t^{-\nu}I_\nu(t)$ is increasing on $(0,\infty)$, for
$0<\sigma\leq\rho$,
\begin{equation}\label{eq_phi_ratio}
    \frac{\phi_{\ell,Y}(\sigma)}{\phi_{\ell,Y}(\rho)}
    =
    \Bigl(\frac{\sigma}{\rho}\Bigr)^{\!-\mathfrak b}
    \frac{I_{\nu_\ell}(\tau_\sigma)}{I_{\nu_\ell}(\tau_\rho)}
    \le
    \Bigl(\frac{\sigma}{\rho}\Bigr)^{\!\gamma_\ell}.
\end{equation}
On the other hand, for $0<\rho\le 2$, there exists $0<c<1$ such that
\begin{equation}\label{eq_psi_comparable_K}
    c\rho^{-\mathfrak b}K_{\nu_\ell}(\tau_\rho)
    \le
    \psi_{\ell,Y}(\rho)
    \le
    \rho^{-\mathfrak b}K_{\nu_\ell}(\tau_\rho),\qquad \ell \ge 0.
\end{equation}
Indeed, as $t^\nu K_\nu(t)$ is decreasing on $(0,\infty)$, the definition of $\eta_{\ell,Y}$ gives
\begin{equation*}
\frac{\eta_{\ell,Y}I_{\nu_\ell}(\tau_\rho)}{K_{\nu_\ell}(\tau_\rho)} =
\frac{I_{\nu_\ell}(\tau_\rho)}{I_{\nu_\ell}(\tau_2)}
\frac{K_{\nu_\ell}(\tau_2)}{K_{\nu_\ell}(\tau_\rho)}\le
\Bigl(\frac{\rho}{2}\Bigr)^{\!2\mathfrak a\nu_\ell}\le
4^{-\mathfrak a\nu_0} < 1.
\end{equation*}
% \begin{equation*}
% \frac{\eta_{\ell,Y}I_{\nu_\ell}(\tau_\rho)}{K_{\nu_\ell}(\tau_\rho)}\le
% \Bigl(\frac{\rho}{2}\Bigr)^{\!2\mathfrak a\nu_\ell}\le
% 4^{-\mathfrak a\nu_0} < 1.
% \end{equation*}
This establishes \eqref{eq_psi_comparable_K}. Consequently, if $\rho\le \sigma \le 2$, then
\begin{equation}\label{eq_psi_ratio}
    \frac{\psi_{\ell,Y}(\sigma)}{\psi_{\ell,Y}(\rho)}
    \leq
    C
    \Bigl(\frac{\rho}{\sigma}\Bigr)^{\!\delta_\ell}.
\end{equation}
By Lemma~\ref{lem_bessel_inputs_high_frequency} and $\nu_\ell \ge \nu_0 >0$,
\begin{equation}\label{eq_IK_product_bound}
I_{\nu_\ell} (t) K_{\nu_\ell} (t) \le \frac{C}{1+\nu_\ell}, \qquad \ell \ge 0, \qquad t>0.
\end{equation}
Combining \eqref{eq_psi_comparable_K} and \eqref{eq_IK_product_bound},
\begin{equation}\label{eq_phi_psi_product_bound}
    \phi_{\ell,Y}(\rho)\psi_{\ell,Y}(\rho)
    \le
    C\rho^{-2\mathfrak b}(1+\ell)^{-1},
    \qquad
    0<\rho\leq1,
    \qquad
    \ell\ge 1.
\end{equation}
Fix $0<\rho\le 1$. Split
\begin{equation*}
\begin{aligned}
    \int_0^2 G_{\ell,Y}(\rho,\sigma) \, \sigma^{n-1} \, d\sigma
    &=
    \mathfrak a^{-1} \psi_{\ell,Y}(\rho)
    \int_0^\rho
        \phi_{\ell,Y}(\sigma)\,\sigma^{n-1}\,d\sigma          \\
    &\quad+
    \mathfrak a^{-1} \phi_{\ell,Y}(\rho)
    \int_\rho^2
        \psi_{\ell,Y}(\sigma)\,\sigma^{n-1}\,d\sigma          \\
    &=: I_1+I_2.
\end{aligned}
\end{equation*}
For $I_1$, \eqref{eq_phi_ratio} together with \eqref{eq_phi_psi_product_bound} implies
\begin{align}\label{eq_I1_bound_schur}
    I_1 &\le \mathfrak a^{-1} \phi_{\ell,Y}(\rho)\psi_{\ell,Y}(\rho) \int_0^\rho \Bigl(\frac{\sigma}{\rho}\Bigr)^{\!\gamma_\ell} \sigma^{n-1} \, d\sigma\\ \nonumber
    &= \frac{\mathfrak a^{-1}}{n+\gamma_\ell} \, \rho^n \phi_{\ell,Y}(\rho)\psi_{\ell,Y}(\rho)\\ \nonumber
    &\le C (1+\ell)^{-2}\rho^{n-2\mathfrak b}\\ \nonumber
    &\le C (1+\ell)^{-2}\rho^{1-\alpha}.
\end{align}
For $I_2$, using \eqref{eq_psi_ratio} and \eqref{eq_phi_psi_product_bound},
\begin{align*}
    I_2 &\le C\phi_{\ell,Y}(\rho)\psi_{\ell,Y}(\rho)
    \int_\rho^2 \Bigl(\frac{\rho}{\sigma}\Bigr)^{\!\delta_\ell}\sigma^{n-1} \, d\sigma \\
    &\le C (1+\ell)^{-1} \rho^{\delta_\ell - 2 \mathfrak b} \int_\rho^2 \sigma^{n-\delta_\ell - 1} \, d\sigma \\
    &\le C \begin{cases}
        (1+\ell)^{-2} \rho^{n-2\mathfrak b}, & \delta_\ell >n,\\
        (1+\ell)^{-1} \rho^{n-2\mathfrak b} |\log{\rho}|, & \delta_\ell = n,\\
        (1+\ell)^{-1} |n-\delta_\ell|^{-1} \rho^{\delta_\ell-n}, &\delta_\ell<n.
    \end{cases}
\end{align*}
For $\delta_\ell \ge n$, it is clear that $I_2 \le C (1+\ell)^{-2} \rho^{1-\alpha}$ since $n-2\mathfrak b=2-2\alpha$ and $\rho\le 1$ (the logarithmic blow-up $|\log{\rho}|$ can be bounded by $\rho^{-\varepsilon}$, with $0<\varepsilon<1-\alpha$). While if $\delta_\ell <n$, first note that for $n\ge 2$, 
\begin{equation}\label{eq_delta_lower}
    \delta_\ell \ge \delta_1\geq n+\alpha-1.
\end{equation}
It follows that $\rho^{\delta_\ell-n} \le \rho^{\delta_1-n} \le \rho^{n+\alpha-1-2\mathfrak b} = \rho^{1-\alpha}$. The coefficients in $\ell$ are harmless as there are only finite many $\ell\ge 1$ such that $\delta_\ell < n$. In summary, 
\begin{equation*}
    I_2 \le C (1+\ell)^{-2} \rho^{1-\alpha},\qquad 0<\rho \le 1.
\end{equation*}
Combining this with \eqref{eq_I1_bound_schur} proves \eqref{eq_basic_R_schur_improved} and \eqref{eq_basic_R_schur} with $\ell \ge 1$.


% To continue, note first that, since $n\ge 2$,
% \begin{equation}\label{eq_delta_lower}
%     \delta_1\geq n+\alpha-1.
% \end{equation}
% A similar computation shows that $\delta_\ell \ge \delta_2 > n$ for $\ell \ge 2$. Hence
% \begin{align}\label{eq_sigma_l>=2}
%     \int_\rho^2 \sigma^{n-\delta_\ell - 1} \, d\sigma \le \frac{\rho^{n-\delta_\ell}}{\delta_\ell - n} \le C(1+\ell)^{-1} \rho^{n-\delta_\ell},\qquad \ell \ge 2.
% \end{align}
% On the other hand, if $\ell = 1$, then
% \begin{align}\label{eq_l=1}
%     \int_\rho^2 \sigma^{n-\delta_1 - 1} \, d\sigma \le \begin{cases}
%         C, & \delta_1\le n,\\[4pt]
%         C \rho^{n-\delta_1}, & \delta_1>n.
%     \end{cases}
% \end{align}
% Therefore, if $\ell \ge 2$ or $\ell=1$ with $\delta_1 >n$, then
% \begin{equation*}
%     I_2 \le C (1+\ell)^{-2} \rho^{n-2\mathfrak b} \le C (1+\ell)^{-2} \rho^{1-\alpha},
% \end{equation*}
% whereas if $\ell =1$ and $\delta_1 \le n$, then
% \begin{equation*}
%     I_2\le C (1+\ell)^{-1} \rho^{\delta_1-2\mathfrak b} \le C \rho^{n+\alpha-1 - 2\mathfrak b} = C \rho^{1-\alpha}.
% \end{equation*}
% Together with \eqref{eq_I1_bound_schur}, this proves \eqref{eq_basic_R_schur_improved}, and \eqref{eq_basic_R_schur} for $\ell \ge 1$ follows at once.

It remains to treat the case $\ell=0$. Observe that \eqref{eq_psi_comparable_K} and \eqref{eq_IK_product_bound} remain valid. From the explicit formula \eqref{eq_kernel_Green_copy} and the monotonicity of $t^{-\nu}I_{\nu}(t)$ and $t^\nu K_\nu(t)$,
\begin{equation*}
    G_{0,Y}(\rho,\sigma) \le C \begin{cases}
        \rho^{\gamma_0} \sigma^{-\delta_0}, & 0<\rho \le \sigma <2,\\[4pt]
        \rho^{-\delta_0} \sigma^{\gamma_0}, & 0<\sigma <\rho<2.
    \end{cases}
\end{equation*}
Therefore, since $\gamma_0 = 0$ and $\delta_0 = n-2(1-\alpha) <n$, for all $0<\rho \le 1$,
\begin{align*}
    \int_0^2 G_{0,Y}(\rho,\sigma)\,\sigma^{n-1}\,d\sigma
    &\le C \rho^{-\delta_0} \int_0^\rho \sigma^{n-1}\,d\sigma\\
    &\quad+ C \int_\rho^2 \sigma^{n-\delta_0 -1}\,d\sigma \\
    &\le C.
\end{align*}
This proves \eqref{eq_basic_R_schur_ell=0} and \eqref{eq_basic_R_schur} with $\ell=0$.




\textit{Case 2: $|Y|> C^*(1+\ell)$.} Set for simplicity $k_{\ell,Y}(\rho):=\rho^{-\mathfrak b}K_{\nu_\ell}(\tau_\rho)$. 
% \begin{equation*}
%     k_{\ell,Y}(\rho):=\rho^{-\mathfrak b}K_{\nu_\ell}(\tau_\rho),\quad 0<\rho <2.
% \end{equation*}
By \eqref{eq_psi_comparable_K}, it is enough to estimate $G_{\ell,Y}(\rho,\sigma)$ with $\psi_{\ell,Y}$ replaced by $k_{\ell,Y}$. Fix $0<\rho<2$. From \eqref{eq_kernel_Green_copy},
\begin{align}\label{eq_weighted_split}
    \rho^{2\beta}
    \int_0^2
        G_{\ell,Y}(\rho,\sigma)
    \, \sigma^{n-1} \, d\sigma                                      \nonumber
    \le
    C\rho^{2\beta}
    \phi_{\ell,Y}(\rho)k_{\ell,Y}(\rho)
    \biggl[\int_0^\rho \frac{\phi_{\ell,Y}(\sigma)}{\phi_{\ell,Y}(\rho)} \, \sigma^{n-1} \, d\sigma + \int_\rho^2 \frac{k_{\ell,Y}(\sigma)}{k_{\ell,Y}(\rho)} \, \sigma^{n-1} \, d\sigma \biggr].
\end{align}
By Lemma~\ref{lem_bessel_inputs_high_frequency}, for $\nu_\ell\ge \nu_0>0$,
\begin{equation}\label{eq_phi_k_product}
    \phi_{\ell,Y}(\rho)k_{\ell,Y}(\rho)
    =
    \rho^{-2\mathfrak b}
    I_{\nu_\ell}(\tau_\rho)K_{\nu_\ell}(\tau_\rho)
    \le
    C\rho^{-2\mathfrak b}(1 + \ell + \tau_\rho)^{-1}.
\end{equation}
Put $r=\sigma/\rho\in(0,1]$. Then
\begin{equation*}
\frac{\phi_{\ell,Y}(\sigma)}{\phi_{\ell,Y}(\rho)} = r^{\gamma_\ell} \frac{\mathcal I_{\nu_\ell}(\tau_\rho r^{\mathfrak a})}{\mathcal I_{\nu_\ell}(\tau_\rho)}.
\end{equation*}
Assume first that $\tau_\rho \le  C_0(1+\ell)$ with $C_0>0$ to be determined later. Since $\mathcal I_\nu$ is increasing,
\begin{align*}
    \int_0^\rho
        \frac{\phi_{\ell,Y}(\sigma)}
             {\phi_{\ell,Y}(\rho)}
    \, \sigma^{n-1} \, d\sigma
    &\le C\rho^n \int_0^1 r^{n-1+\gamma_\ell}\,dr
    \le C\frac{\rho^n}{1+\ell}
    \le C\frac{\rho^n}{1+ \ell + \tau_\rho}.
\end{align*}
Next assume that $\tau_\rho>C_0(1+\ell)$. Choose $C_0>1$ sufficiently large so that $2^{- \mathfrak a} C_0 (1+\ell) \ge C(1+\ell)$, where $C>0$ is the constant from \eqref{eq_S_lower_large}. Lemma~\ref{lem_bessel_inputs_high_frequency} \eqref{eq_S_lower_large} then gives for $0<r<1/2$ (thus $r^{\mathfrak a}\le 2^{- \mathfrak a}$),
\begin{align*}
\frac{\mathcal I_{\nu_\ell}(\tau_\rho r^{\mathfrak a})}{\mathcal I_{\nu_\ell}(\tau_\rho)}
&=
\exp\Bigl\{-\int_{\tau_\rho r^{\mathfrak a}}^{\tau_\rho}\frac{I_{\nu_\ell+1}(t)}{I_{\nu_\ell}(t)}\,dt\Bigr\} \\
&\le
\exp\Bigl\{-\int_{\tau_\rho 2^{-\mathfrak a}}^{\tau_\rho}\frac{I_{\nu_\ell+1}(t)}{I_{\nu_\ell}(t)}\,dt\Bigr\}
\le C e^{-c \tau_\rho},
\end{align*}
and for $1/2\le r \le 1$ (thus $\tau_\rho r^{\mathfrak a} \ge 2^{- \mathfrak a} C_0 (1+\ell) \ge C(1+\ell)$),
\begin{equation*}
\frac{\mathcal I_{\nu_\ell}(\tau_\rho r^{\mathfrak a})}{\mathcal I_{\nu_\ell}(\tau_\rho)}
=\exp\Bigl\{-\int_{\tau_\rho r^{\mathfrak a}}^{\tau_\rho}\frac{I_{\nu_\ell+1}(t)}{I_{\nu_\ell}(t)}\,dt\Bigr\}
\le C e^{-c \tau_\rho(1-r^{\mathfrak a})}
\le C e^{-c \tau_\rho(1-r)}.
\end{equation*}
Hence
\begin{align*}
   \int_0^\rho \frac{\phi_{\ell,Y}(\sigma)}{\phi_{\ell,Y}(\rho)}\,\sigma^{n-1}\,d\sigma
   &\le
    C\rho^n \int_0^{1/2} r^{n-1+\gamma_\ell} e^{-c\tau_\rho}\,dr\\
    &\quad+ C\rho^n \int_{1/2}^1 r^{n-1+\gamma_\ell}e^{-c\tau_\rho(1-r)}\,dr\\
    &\le C\frac{\rho^n}{1+\ell+\tau_\rho}.
\end{align*}
Thus, in all cases,
\begin{equation}\label{eq_I_ratio_integral_final}
    \int_0^\rho
        \frac{\phi_{\ell,Y}(\sigma)}
             {\phi_{\ell,Y}(\rho)}
    \, \sigma^{n-1} \, d\sigma
    \le
    C\frac{\rho^n}{1+\ell+\tau_\rho}.
\end{equation}

For the $K$-branch, note that now $r=\sigma/\rho \ge 1$. Then
\begin{equation}\label{eq_K_ratio}
    \frac{k_{\ell,Y}(\sigma)}
         {k_{\ell,Y}(\rho)}
    =
    r^{-\delta_\ell}
    \frac{\mathcal K_{\nu_\ell}(\tau_\rho r^{\mathfrak a})}
         {\mathcal K_{\nu_\ell}(\tau_\rho)},
    \qquad
    \mathcal K_\nu(t):=t^\nu K_\nu(t),
\end{equation}
where
\begin{equation*}
    \delta_\ell=\mathfrak b+\mathfrak a\nu_\ell =
    \mathfrak b+\sqrt{\mathfrak b^2+\lambda_\ell}
    \sim 1+\ell.
\end{equation*}
Moreover, for $\ell\ge 1$, \eqref{eq_delta_lower} implies $\delta_\ell\ge n+\alpha-1$.

\textit{Case 2a: $\delta_\ell > n$.} If $\tau_\rho\le C_0(1+\ell)$, then $\mathcal K_\nu$ is decreasing and hence
\begin{align}\label{eq_formula3}
\int_\rho^2 \frac{k_{\ell,Y}(\sigma)}{k_{\ell,Y}(\rho)}\,\sigma^{n-1}\,d\sigma
&\le \rho^n
    \int_1^{2/\rho} r^{n-1-\delta_\ell}\,dr
    \le C\frac{\rho^n}{1+\ell}
    \le C\frac{\rho^n}{1+ \ell + \tau_\rho}.
\end{align}
If instead $\tau_\rho>C_0 (1+\ell)$, where $C_0$ is again chosen large enough so that Lemma~\ref{lem_bessel_inputs_high_frequency} applies, then \eqref{eq_K_normalized_exp_decay} yields
\begin{equation*}
    \frac{\mathcal K_{\nu_\ell}(\tau_\rho r^{\mathfrak a})}
         {\mathcal K_{\nu_\ell}(\tau_\rho)}
    \le
    C e^{-c\tau_\rho(r^{\mathfrak a}-1)}
    \le C e^{-c\tau_\rho(r-1)}
\end{equation*}
since $r\ge 1$. Therefore,
\begin{equation}\label{eq_K_ratio_large_ell}
    \int_\rho^2\frac{k_{\ell,Y}(\sigma)}{k_{\ell,Y}(\rho)}\,\sigma^{n-1}\,d\sigma
    \le C\rho^n  \int_1^\infty e^{-c\tau_\rho(r-1)}\,dr
    \le C\frac{\rho^n}{1+\ell+\tau_\rho}.
\end{equation}
Combining this with \eqref{eq_phi_k_product},
\begin{align*}
    \rho^{2\beta}
    \int_0^2
        G_{\ell,Y}(\rho,\sigma)
    \, \sigma^{n-1} \, d\sigma
    \le \frac{C \rho^{n-2 \mathfrak b+2\beta} }{(1+\ell + \tau_\rho)^2}
    = \frac{C \rho^{2 \mathfrak a} }{(1+\ell + \tau_\rho)^2}
    \le C |Y|^{-2},
\end{align*}
as $\rho^{2 \mathfrak a} = \mathfrak a^2 |Y|^{-2} \tau_\rho^2$. This proves \eqref{eq_basic_R_schur} with $\delta_\ell >n$.

\textit{Case 2b: $\delta_\ell \le n$.} On the region $\sigma<\rho$, the change of variables $u=\tau_\sigma$ gives
\begin{align}\label{eq_finite_left_scaling}
\rho^{2\beta}
    k_{\ell,Y}(\rho)
    \int_0^\rho
        \phi_{\ell,Y}(\sigma) \, \sigma^{n-1} \, d\sigma
    =
    C |Y|^{-2}
    \tau_\rho^{\frac{2\beta-\mathfrak b}{\mathfrak a}}
    K_{\nu_\ell}(\tau_\rho)
    \int_0^{\tau_\rho}
        u^{\frac{n-\mathfrak b}{\mathfrak a}-1}
        I_{\nu_\ell}(u) \, du.
\end{align}
Since only finitely many values of $\ell$ are involved here, the standard small and large argument asymptotic formulas for the modified Bessel functions suffice, and there is no need to keep track of the $\ell$-dependence. We show that the function of $\tau_\rho$ on the right-hand side of \eqref{eq_finite_left_scaling} is uniformly bounded on $(0,\infty)$. By Lemma~\ref{lem:Bessel:integral} $(ii)$, this function is nothing but $II(q,s,\nu_\ell,\tau_\rho)$ with $q=\frac{2\beta-\mathfrak b}{\mathfrak a}$, $s=\frac{n-\mathfrak b}{\mathfrak a}$ and $\nu_\ell \ge \nu_0>0$. It is easy to check 
\begin{align}\label{check}
    q+s+\nu-|\nu| = q+s = \frac{2\beta - 2\mathfrak b+n}{\mathfrak a} = 2 \le 2,
\end{align}
inferring the uniform boundedness of $II(q,s,\nu_\ell,\tau_\rho)$ for all $\tau_\rho >0$, and hence \eqref{eq_finite_left_scaling} is $O(|Y|^{-2})$ as required.

On the other hand,
% Near $\tau_\rho=0$,
% \begin{equation*}
%     \tau_\rho^{\frac{2\beta-\mathfrak b}{\mathfrak a}}
%     K_{\nu_\ell}(\tau_\rho)
%     \int_0^{\tau_\rho}
%         u^{\frac{n-\mathfrak b}{\mathfrak a}-1}
%         I_{\nu_\ell}(u) \, du
%     \sim \tau_\rho^{\frac{2\beta-\mathfrak b}{\mathfrak a} - \nu_\ell} \int_0^{\tau_\rho} u^{\frac{n-\mathfrak b}{\mathfrak a}-1+\nu_\ell}\,du
%     \sim \tau_\rho^2 \le 1.
% \end{equation*}
% As $\tau_\rho \to \infty$, the contribution coming from $\int_0^1 u^{\frac{n-\mathfrak b}{\mathfrak a}-1} I_{\nu_\ell}(u) \, du$
% % \begin{equation*}
% %     \int_0^1 u^{\frac{n-\mathfrak b}{\mathfrak a}-1} I_{\nu_\ell}(u) \, du
% % \end{equation*}
% is uniformly bounded because of the exponential decay of $K_{\nu_\ell}(\tau_\rho)$. For the remaining part,
% \begin{align}\label{eq_formula1}
%     \tau_\rho^{\frac{2\beta-\mathfrak b}{\mathfrak a}}
%     K_{\nu_\ell}(\tau_\rho)
%     \int_1^{\tau_\rho}
%         u^{\frac{n-\mathfrak b}{\mathfrak a}-1}
%         I_{\nu_\ell}(u) \, du
%     \sim \tau_\rho^{\frac{2\beta-\mathfrak b}{\mathfrak a}-\frac{1}{2}} e^{-\tau_\rho} \int_1^{\tau_\rho} u^{\frac{n-\mathfrak b}{\mathfrak a}-\frac{3}{2}} e^{u} \, du.
% \end{align}
% If $\frac{n-\mathfrak b}{\mathfrak a}\ge \frac{3}{2}$, then \eqref{eq_formula1} is $O\bigl( \tau_\rho^{\frac{2\beta-\mathfrak b}{\mathfrak a} -\frac{1}{2}+\frac{n-\mathfrak b}{\mathfrak a} - \frac{3}{2}} \bigr) = O(1)$.
% % \begin{equation*}
% %     O\bigl( \tau_\rho^{\frac{2\beta-\mathfrak b}{\mathfrak a} -\frac{1}{2}+\frac{n-\mathfrak b}{\mathfrak a} - \frac{3}{2}} \bigr) = O(1).
% % \end{equation*}
% While if $\frac{n-\mathfrak b}{\mathfrak a} < \frac{3}{2}$, then
% \begin{align*}
% \int_1^{\tau_\rho} u^{\frac{n-\mathfrak b}{\mathfrak a}-\frac{3}{2}} e^{-(\tau_\rho-u)}\,du
% &\le e^{-\tau_\rho/2} \int_1^{\tau_\rho/2} u^{\frac{n-\mathfrak b}{\mathfrak a}-\frac{3}{2}}\,du\\
% &\quad+ C \tau_\rho^{\frac{n-\mathfrak b}{\mathfrak a}-\frac{3}{2}} \int_{\tau_\rho/2}^{\tau_\rho} e^{-(\tau_\rho - u)}\,du\\
% &\le C \tau_\rho^{\frac{n-\mathfrak b}{\mathfrak a}-\frac{3}{2}},
% \end{align*}
% and \eqref{eq_formula1} is again $O(1)$. Hence, for $\delta_\ell \le n$, \eqref{eq_finite_left_scaling} is bounded by $C |Y|^{-2}$.
consider the region $\sigma \ge \rho$. After the same change of variables,
\begin{align*}
\rho^{2\beta}
    \phi_{\ell,Y}(\rho)
    \int_\rho^2
        k_{\ell,Y}(\sigma)\,\sigma^{n-1}\,d\sigma
    =
    C |Y|^{-2}
    \tau_\rho^{\frac{2\beta-\mathfrak b}{\mathfrak a}}
    I_{\nu_\ell}(\tau_\rho)
    \int_{\tau_\rho}^\infty
        u^{\frac{n-\mathfrak b}{\mathfrak a}-1}
        K_{\nu_\ell}(u) \, du = I(q,s,\nu_\ell
        ,0,\tau_\rho),
\end{align*}
with the same $q,s$ as in the case $\sigma>\rho$. With \eqref{check} and 
\begin{align*}
q+\nu_\ell = \frac{2\beta-\mathfrak b + \sqrt{\mathfrak b^2 + \lambda_\ell}}{\mathfrak a} \ge \frac{2\beta}{\mathfrak a}\ge 0,
\end{align*}
Lemma~\ref{lem:Bessel:integral} gives $I(\tau_\rho) \le C$ for all $\tau_\rho >0$. This closes the proof of \eqref{eq_basic_R_schur} for $\delta_\ell\le n$.

% bounded on $(0,\infty)$. Near $\tau_\rho=0$,
% \begin{align*}
%     \tau_\rho^{\frac{2\beta-\mathfrak b}{\mathfrak a}}
%     I_{\nu_\ell}(\tau_\rho)
%     \int_{\tau_\rho}^\infty
%         u^{\frac{n-\mathfrak b}{\mathfrak a}-1}
%         K_{\nu_\ell}(u) \, du
%     &=
%     \tau_\rho^{\frac{2\beta-\mathfrak b}{\mathfrak a}}
%     I_{\nu_\ell}(\tau_\rho)
%     \int_{\tau_\rho}^1
%         u^{\frac{n-\mathfrak b}{\mathfrak a}-1}
%         K_{\nu_\ell}(u) \, du\\
%     &\quad+ \tau_\rho^{\frac{2\beta-\mathfrak b}{\mathfrak a}}
%     I_{\nu_\ell}(\tau_\rho)
%     \int_{1}^\infty
%         u^{\frac{n-\mathfrak b}{\mathfrak a}-1}
%         K_{\nu_\ell}(u) \, du\\
%     &\le C \tau_\rho^{\frac{2\beta-\mathfrak b}{\mathfrak a}+\nu_\ell} \int_{\tau_\rho}^1 u^{\frac{n-\mathfrak b}{\mathfrak a}-\nu_\ell-1}\,du \\
%     &\quad+ C \tau_\rho^{\frac{2\beta-\mathfrak b}{\mathfrak a}+\nu_\ell} \int_1^\infty u^{\frac{n-\mathfrak b}{\mathfrak a}-\frac{3}{2}} e^{-u}\,du.
% \end{align*}
% The last term is plainly $O\bigl( \tau_\rho^{\frac{2\beta-\mathfrak b}{\mathfrak a}+\nu_\ell} \bigr) = O(1)$. For the preceding term, if $\frac{n-\mathfrak b}{\mathfrak a}-\nu_\ell <0$, then it is $O\bigl( \tau_\rho^{\frac{2\beta-\mathfrak b}{\mathfrak a}+\nu_\ell + \frac{n-\mathfrak b}{\mathfrak a}-\nu_\ell} \bigr) = O(1)$. If instead $\frac{n-\mathfrak b}{\mathfrak a}-\nu_\ell \ge 0$, then it is bounded by $O\bigl( \tau_\rho^{\frac{2\beta-\mathfrak b}{\mathfrak a}+\nu_\ell} |\log \tau_\rho| \bigr) = O(1)$, as $\tau_\rho \to 0$. Now suppose $\tau_\rho \to \infty$. Then
% \begin{align}\label{eq_formula2}
%     \tau_\rho^{\frac{2\beta-\mathfrak b}{\mathfrak a}}
%     I_{\nu_\ell}(\tau_\rho)
%     \int_{\tau_\rho}^\infty
%         u^{\frac{n-\mathfrak b}{\mathfrak a}-1}
%         K_{\nu_\ell}(u) \, du
%     \le C \tau_\rho^{\frac{2\beta-\mathfrak b}{\mathfrak a}-\frac{1}{2}} e^{\tau_\rho} \int_{\tau_\rho}^\infty u^{\frac{n-\mathfrak b}{\mathfrak a} - \frac{3}{2}} e^{-u}\,du.
% \end{align}
% If $\frac{n-\mathfrak b}{\mathfrak a} - \frac{3}{2} \le 0$, then the right-hand side is bounded by
% \begin{equation*}
%     C \tau_\rho^{\frac{2\beta-\mathfrak b}{\mathfrak a}-\frac{1}{2} + \frac{n-\mathfrak b}{\mathfrak a} - \frac{3}{2} } e^{\tau_\rho} \int_{\tau_\rho}^\infty e^{-u}\,du \le C.
% \end{equation*}
% If $\frac{n-\mathfrak b}{\mathfrak a} - \frac{3}{2} > 0$, then using $\tau_\rho + u \le \tau_\rho(1+u)$, \eqref{eq_formula2} is bounded by
% \begin{align*}
%     C \tau_\rho^{\frac{2\beta-\mathfrak b}{\mathfrak a}-\frac{1}{2}}\int_{\tau_\rho}^\infty u^{\frac{n-\mathfrak b}{\mathfrak a} - \frac{3}{2}} e^{-(u-\tau_\rho)}\,du
%     &= C \tau_\rho^{\frac{2\beta-\mathfrak b}{\mathfrak a}-\frac{1}{2}}\int_{0}^\infty (\tau_\rho + u)^{\frac{n-\mathfrak b}{\mathfrak a} - \frac{3}{2}} e^{-u}\,du \\
%     &\le C \int_0^\infty (1+u)^{\frac{n-\mathfrak b}{\mathfrak a} - \frac{3}{2}} e^{-u}\,du \le C.
% \end{align*}
% Thus \eqref{eq_basic_R_schur} also holds for $\delta_\ell \le n$.

To complete the argument, it remains to derive a pointwise estimate for the $K$-branch when $\delta_\ell \le n$. Clearly, \eqref{eq_K_ratio_large_ell} remains valid when $\tau_\rho > C_0(1+\ell)$. For the case $\tau_\rho \le C_0(1+\ell)$, however, \eqref{eq_delta_lower} only gives
\begin{align*}
    \int_\rho^2 \frac{k_{\ell,Y}(\sigma)}{k_{\ell,Y}(\rho)} \, \sigma^{n-1} \, d\sigma
    &\le C \rho^n \int_1^{2/\rho} r^{n-\delta_\ell-1}\,dr\\
    &\le C \begin{cases}
        \rho^n |\log{\rho}|, & \delta_\ell = n,\\
        |n-\delta_\ell|^{-1} \rho^{\delta_\ell}, & \delta_\ell < n.
    \end{cases}
\end{align*}
For $\delta_\ell=n$, it is plain that the integral is bounded by $C(1+\ell+\tau_\rho)^{-1} \rho^{n+\alpha-1}$, where one can bound the logarithmic term by $C_\alpha \rho^{-(1-\alpha)}$. As for $\delta_\ell<n$, it satisfies upper bound
\begin{align}\label{equation2}
    \begin{cases}
        (1+\ell+\tau_\rho)^{-1} \rho^{n+\alpha-1}, & \ell \ge 1,\\
        (1+\tau_\rho)^{-1} \rho^{2\mathfrak b}, & \ell=0.
    \end{cases}
\end{align}
% \begin{align*}
%     \int_\rho^2 \frac{k_{\ell,Y}(\sigma)}{k_{\ell,Y}(\rho)} \, \sigma^{n-1} \, d\sigma
%     &\le C \rho^n \int_1^{2/\rho} r^{n-\delta_\ell-1}\,dr\\
%     &\le C \rho^{n} \int_1^{2/\rho} r^{1-\alpha-1}\,dr\\
%     &\le C \rho^{n+\alpha-1}
%     \le C \frac{\rho^{n+\alpha-1}}{1+\ell+\tau_\rho}
% \end{align*}
% for $\ell = 1$, and it is $O((1+\tau_\rho)^{-1}\rho^{\delta_0}) = O((1+\tau_\rho)^{-1} \rho^{2\mathfrak b})$ for $\ell = 0$, since $\delta_0=2\mathfrak b = n+2\alpha-2<n$.
Concluded from \eqref{eq_phi_k_product}, \eqref{eq_I_ratio_integral_final}, \eqref{eq_K_ratio_large_ell} and \eqref{equation2}, for all $0<\rho \le 1$,
\begin{align*}
    \int_0^2
        G_{\ell,Y}(\rho,\sigma)
    \, \sigma^{n-1} \, d\sigma
    &\le C \rho^{-2 \mathfrak b} (1+\ell+\tau_\rho)^{-1} \begin{cases}
        (1+\ell+\tau_\rho)^{-1} \rho^{n+\alpha-1}, & \ell \ge 1,\\
        (1+\tau_\rho)^{-1} \rho^{2\mathfrak b}, & \ell=0.
    \end{cases} \\
    &\le C \begin{cases}
        (1+\ell+\tau_\rho)^{-2} \rho^{1-\alpha}, & \ell \ge 1,\\
        (1+\tau_\rho)^{-2}, & \ell=0.
    \end{cases} 
\end{align*}
% for $\ell = 1$, while for $\ell=0$ it is $(1+\tau_\rho)^{-1} O(\rho^{-2\mathfrak b} \rho^{2 \mathfrak b}) = (1+\tau_\rho)^{-1} O(1)$. 
This completes the proof of \eqref{eq_basic_R_schur_improved}, \eqref{eq_basic_R_schur_ell=0}, and hence Proposition~\ref{prop:Green:Shur}.

\end{proof}

\end{document}
